%% ma: L12-B19-0009
%% lop: 12
%% chuong: 6
%% bai: 19
%% dang: 12.19.4
%% loai: TLN
%% muc: VD
%% dap_an: "0,3"
%% nguon: "(NBV)-ĐỀ SỐ 5-MỖI NGÀY 1 ĐỀ THI 2025.pdf (D05, MỖI NGÀY 1 ĐỀ - Đề số 5), Phần III Câu 5"
%% kien_thuc: "Bài 19 (công thức xác suất toàn phần); Bài 18 (xác suất có điều kiện)"
%% trang_thai: cho-duyet
%% ly_do: "Dạng toán mới đề xuất 12.19.4: xác suất ở lần rút thứ hai khi rút không hoàn lại (công thức xác suất toàn phần)"
%% ghi_chu: "Bỏ ảnh người rút thẻ: ảnh minh họa đời sống, không chứa dữ kiện toán"
\begin{ex}
Có $40$ tấm thẻ kích thước như nhau và đánh số thứ tự lần lượt từ $1$ đến $40$ (mỗi tấm thẻ chỉ ghi một số nguyên dương, hai thẻ khác nhau ghi hai số khác nhau). Một người lần lượt rút hai thẻ (rút không hoàn lại). Tính xác suất lần thứ hai rút được thẻ ghi số nguyên tố.
\shortans{0,3}
\loigiai{Từ $1$ đến $40$ có $12$ số nguyên tố ($2,3,5,7,11,13,17,19,23,29,31,37$). Gọi $A$: ``lần thứ nhất rút được thẻ ghi số nguyên tố'', $B$: ``lần thứ hai rút được thẻ ghi số nguyên tố''.
$P(A)=\dfrac{12}{40}$, $P(\overline A)=\dfrac{28}{40}$, $P(B\mid A)=\dfrac{11}{39}$, $P(B\mid\overline A)=\dfrac{12}{39}$.
Theo công thức xác suất toàn phần: $P(B)=\dfrac{12}{40}\cdot\dfrac{11}{39}+\dfrac{28}{40}\cdot\dfrac{12}{39}=\dfrac{468}{1560}=0{,}3$.}
\end{ex}
